\section{Evaluation}

\subsection{Setup and Methodology}
Our evaluation is designed to avoid ``numbers in a vacuum'' by tying every reported metric to a concrete artifact produced by the toolchain and checked by scripts. For each benchmark, we (i) run the deterministic tick semantics for a fixed number of ticks, (ii) check bit-exactness against a golden digest when a verification manifest is available, and (iii) collect boardless hardware evidence from HLS and (when available) Vivado implementation reports.

The evaluation tables are generated from per-benchmark \texttt{bench\_report.json} outputs and contain: (a) correctness/bit-exact status (\texttt{results\_bitexact.csv}), (b) HLS/Vivado QoR summaries (\texttt{results\_qor.csv}), and (c) a throughput-oriented view that combines cycle-level results with an explicitly stated estimate model (\texttt{results\_throughput.csv}). All quantitative values shown in this section come from those generated tables (via \texttt{\input}), not from manually hardcoded numbers in the manuscript.

\subsection{Correctness / Bit-Exact}
Table~\ref{tab:bitexact-summary} reports the script-generated bit-exact summary used by this paper, including graph size $(N,E)$, tick count, verifier outcome, mismatch count, digest-match status, independent-check status, and trace presence.

\paragraph{Verified bit-exactness.}
For B1, B2, B3, and B4, Table~\ref{tab:bitexact-summary} reports \texttt{verify\_ok=true}, \texttt{mismatches\_len=0}, and \texttt{digestMatchOk=true}. This is the strongest correctness claim in this paper: for manifest-verified benchmarks, the reference and compiled executions agree bit-for-bit on the recorded digest.

\paragraph{Unverified-but-digested runs.}
For B5 and B6, Table~\ref{tab:bitexact-summary} reports \texttt{digestMatchOk=true} and \texttt{trace\_present=true}, while \texttt{verify\_ok} and \texttt{mismatches\_len} remain ``-''. This indicates digested outputs are present, but manifest-style mismatch counting is not asserted by this table for those rows.

\paragraph{Independent-check coverage.}
Table~\ref{tab:bitexact-summary} explicitly separates independent-check coverage from manifest verification by reporting \texttt{independent\_check\_ok} per benchmark row.

\begin{table*}[t]
\centering
\scriptsize
\begin{tabular}{|l|l|c|c|c|c|l|l|}
\hline
benchmarkId & modelId & verify\_ok & mismatches\_len & digestMatchOk & ticks & ir\_sha256 & bench\_report\_path \\
\hline
B1 & example\_b1\_small\_subgraph & true & 0 & true & 20 & 31b5a208c287eaec92e2e39f4f19442a7c080a2ff000c0f667cd984b7e1e8286 & build/paperA\_routeA/B1/example\_b1\_small\_subgraph/bench\_report.json \\
B2 & B2\_mid\_64\_1024 & true & 0 & true & 20 & 61d1c9b1b97d3d8509c59349445472ed8cf76bf0205e560a1cacffd21509e98b & build/paperA\_routeA/B2/B2\_mid\_64\_1024/bench\_report.json \\
B3 & B3\_kernel\_302\_7500 & true & 0 & true & 20 & 05ce6de59b1518129c3d690c5fca81e507dd8c109193fd83b2ae1301e8447080 & build/paperA\_routeA/B3/B3\_kernel\_302\_7500/bench\_report.json \\
B4 & B4\_celegans\_external\_bundle & true & 0 & true & 20 & 5b21d1db4975db16c02802eaba409f4ac194535880a15b215f7bd02d01071473 & build/paperA\_routeA/B4/B4\_celegans\_external\_bundle/bench\_report.json \\
B6 & B6\_delay\_small & - & - & true & 20 & ab9f80fb981df0f73b1bcd925a962c71feece6fb4d0f731479024997be935c34 & build/paperA\_routeA/B6/B6\_delay\_small/bench\_report.json \\
\hline
\end{tabular}

\caption{Bit-exact summary generated from artifact evidence (source: \texttt{papers/paperA/artifacts/tables/results\_bitexact.csv}).}
\label{tab:bitexact-summary}
\end{table*}

\subsection{Hardware Evidence (boardless)}
Table~\ref{tab:qor-summary} summarizes boardless hardware evidence from generated reports, including HLS schedule fields (\texttt{ii}, \texttt{latencyCycles}), resource utilization fields (\texttt{lut}, \texttt{ff}, \texttt{bram}, \texttt{dsp}), and implementation/timing fields when available.

\paragraph{What \texttt{ii} and \texttt{latencyCycles} mean.}
We interpret each \emph{tick} as one full deterministic network update step. In this context:
\begin{itemize}
\item \textbf{Initiation interval (II):} \texttt{ii} is the cycle spacing between starts of consecutive ticks in steady state.
\item \textbf{Latency cycles:} \texttt{latencyCycles} is the pipeline latency of one tick.
\end{itemize}
Thus, II controls steady-state tick throughput, while \texttt{latencyCycles} controls fill/flush transients.

\paragraph{HLS evidence.}
All listed benchmarks in Table~\ref{tab:qor-summary} report \texttt{hls\_csim\_ok=true} and \texttt{hls\_csynth\_ok=true}, showing that HLS C-simulation and HLS synthesis completed for each listed row.

\paragraph{Vivado implementation coverage is not universal.}
Table~\ref{tab:qor-summary} shows \texttt{vivado\_impl\_ok=true} for B1/B2/B3 and \texttt{vivado\_impl\_ok=false} for B4/B5/B6, with missing timing fields represented explicitly (for example \texttt{wns='-'} on non-implemented rows). Accordingly, the strongest boardless ``scale + implementation'' evidence in this artifact pack is B3, which is both the largest listed benchmark and marked with successful implementation/timing fields.

\paragraph{Resource context.}
Utilization and timing fields are tool-reported boardless evidence from the flow outputs; they are not on-board measurements.

\begin{table*}[t]
\centering
\scriptsize
\begin{tabular}{|l|c|c|c|c|c|c|c|c|c|c|c|c|l|}
\hline
benchmarkId & ii & latencyCycles & lut & ff & bram & dsp & wns & csim\_ok & csynth\_ok & cosim\_ok & vivado\_impl\_ok & vivado\_impl\_status & vivado\_impl\_reason \\
\hline
B1 & 32 & 31 & 1148 & 355 & 0 & 0 & 1.129 & true & true & - & true & OK & - \\
B2 & 614 & 613 & 4818 & 4738 & 1 & 1 & 0.957 & true & true & - & true & OK & - \\
B3 & 2746 & 2745 & 3172 & 3110 & 7 & 1 & 0.988 & true & true & - & true & OK & - \\
B4 & 100 & 99 & 2584 & 2633 & 0 & 0 & 1.032 & true & true & - & true & OK & - \\
B6 & 67 & 66 & 1928 & 1098 & 1 & 0 & 0.894 & true & true & - & true & OK & - \\
\hline
\end{tabular}

\caption{QoR summary (boardless) generated from \texttt{bench\_report.json} artifacts (source: \texttt{papers/paperA/artifacts/tables/results\_qor.csv}).}
\label{tab:qor-summary}
\end{table*}

\subsection{Throughput Context (boardless estimate + CPU baseline)}
Table~\ref{tab:throughput-summary} contextualizes QoR by expressing throughput in ticks/s and pairing estimated HW throughput with a measured CPU baseline.

\paragraph{Cycle-to-throughput translation.}
For estimated maximum frequency \texttt{fmax\_est} (MHz) and initiation interval \texttt{ii} (cycles/tick), steady-state estimated HW tick throughput is:
\[
\texttt{hw\_ticks\_per\_sec\_est} = \frac{\texttt{fmax\_est} \cdot 10^{6}}{\texttt{ii}}.
\]
This gives the same unit (ticks/s) as the CPU baseline and makes II/frequency comparisons interpretable.

\paragraph{Estimated HW vs measured CPU.}
Table~\ref{tab:throughput-summary} reports both \texttt{cpp\_ticks\_per\_sec} and \texttt{hw\_ticks\_per\_sec\_est} per benchmark row. Benchmarks marked with failed implementation coverage in Table~\ref{tab:qor-summary} remain explicitly flagged in throughput notes (for example \texttt{vivado=FAIL}), so those estimates are not interpreted as timing-closure evidence.

\begin{table*}[t]
\centering
\scriptsize
\begin{tabular}{|l|r|r|r|r|r|l|}
\hline
benchmarkId & CPU ticks/s & II (cyc/tick) & Fmax est. (MHz) & HW est. ticks/s & speedup est. & vivado impl status \\
\hline
B1 & 495514.358427 & 32.000000 & 258.331181 & 8072849.392922 & 16.291858 & OK \\
B2 & 162602.683651 & 614.000000 & 247.341083 & 402835.640642 & 2.477423 & OK \\
B3 & 44480.467700 & 2746.000000 & 249.252243 & 90769.207309 & 2.040653 & OK \\
\hline
\end{tabular}

\caption{Throughput context from generated artifacts: measured CPU ticks/s and boardless HW ticks/s estimates (source: \texttt{papers/paperA/artifacts/tables/results\_throughput.csv}).}
\label{tab:throughput-summary}
\end{table*}

\subsection{Benchmark Coverage and Interpretation}
The evaluated set spans correctness, scaling, external-graph flow, and delay-aware variants, as reflected directly in the generated rows for B1/B2/B3/B4/B5/B6 in Tables~\ref{tab:bitexact-summary}--\ref{tab:throughput-summary}. The coverage artifact \texttt{papers/paperA/artifacts/evidence/vivado\_coverage\_report.md} is included to keep implementation coverage auditable and consistent with \texttt{vivado\_impl\_ok} in Table~\ref{tab:qor-summary}.

\subsection{What this does \emph{not} prove}
\begin{itemize}
\item These results are not on-board performance or power measurements; hardware evidence here is tool-based.
\item For rows with \texttt{vivado\_impl\_ok=false}, reported \texttt{hw\_ticks\_per\_sec\_est} does not prove timing closure or deployability.
\item The CPU baseline is a minimal context anchor and not a claim of highly tuned CPU optimization.
\item QoR/utilization fields are tool-reported and do not replace on-board validation.
\end{itemize}

\subsection{Roadmap to on-board measurement (future work)}
Once a physical FPGA board is available, the artifact pack can be extended while preserving the same script-verifiable contract:
\begin{itemize}
\item \textbf{Measure real tick throughput:} add on-device ticks/s outputs with raw logs.
\item \textbf{Measure power/energy:} add board telemetry artifacts and raw measurement logs.
\item \textbf{Validate on-board timing closure context:} record exact part, clocking, bitstream, and implementation logs.
\item \textbf{Cross-check on-board correctness:} compare on-device digest logs against golden digests.
\end{itemize}

\paragraph{Claim-to-table pointers (Evaluation).}
\begin{itemize}
\item \textbf{E1 (Bit-exactness for manifest-verified benchmarks):} Table~\ref{tab:bitexact-summary} from \texttt{results\_bitexact.csv} (fields \texttt{verify\_ok}, \texttt{mismatches\_len}, \texttt{digestMatchOk}).
\item \textbf{E2 (Boardless HLS success + explicit Vivado coverage):} Table~\ref{tab:qor-summary} from \texttt{results\_qor.csv} (fields \texttt{hls\_csim\_ok}, \texttt{hls\_csynth\_ok}, \texttt{vivado\_impl\_ok}, \texttt{wns}, \texttt{fmax\_est}).
\item \textbf{E3 (Throughput contextualization):} Table~\ref{tab:throughput-summary} from \texttt{results\_throughput.csv} (fields \texttt{cpp\_ticks\_per\_sec}, \texttt{hw\_ticks\_per\_sec\_est}, \texttt{notes}).
\item \textbf{E4 (Implementation coverage evidence):} \texttt{papers/paperA/artifacts/evidence/vivado\_coverage\_report.md}, consistent with Table~\ref{tab:qor-summary}.
\end{itemize}
